\section{Fault-localization protocol}
\label{ssec:faults}

The evaluator's claim to localize failures is itself validated, by
controlled fault injection under masking.

\paragraph{The artifact chain.} Each evaluated arm is reconstructed as a
chain of concrete artifacts:
\begin{center}\footnotesize
  evaluation\_key $+$ entity\_artifact $\;\rightarrow\;$ relation\_artifact
  $\;\rightarrow\;$ serialized\_graph $\;\rightarrow\;$ grounded\_candidates
  $\;\rightarrow\;$ answer
\end{center}
Every stage checker consumes a declared artifact and nothing downstream of
it, so each artifact has exactly one stage that owns its validity.

\paragraph{The injection protocol.} Each injection mutates \textbf{exactly
one} artifact in the chain and leaves every other artifact untouched. Eight
fault classes are injected, each crossed with three relation types
(\texttt{NEAR}, \texttt{ON\_ENTITY\_SURFACE}, \texttt{ATTACHED\_TO}), giving
24 injected runs, plus clean-artifact runs as controls. The expected
localization of a fault class is the first gating stage whose independent
checker consumes the mutated artifact --- declared per class before any
attribution is read.

\paragraph{The masking protocol.} The evaluator is never told the injected
class. It attributes each failure via its independent per-stage checkers,
exactly as it would on real data, and the attribution is compared to the
hidden injection label only afterward. A localization counts only if the
attributed stage matches the declared expectation for the hidden class.

\paragraph{Result.} 24/24 controlled artifact-level injected faults
localized while the evaluator was masked to the injected class, with zero
spurious failures on clean artifacts.

%% INTEGRATOR: verify this table against Agent B's committed
%% fault-injection results before submission -- the class names, mutated
%% artifacts and expected stages below are placeholders and MUST be replaced
%% by the committed declarations, not paraphrased from memory.
\begin{table}[t]
  \centering\footnotesize
  \setlength{\tabcolsep}{2pt}
  \begin{tabular}{@{}lllc@{}}
    \toprule
    fault class & mutated artifact & expected stage & loc.\\
    \midrule
    \TODO{class 1} & \TODO{artifact} & \TODO{stage} & \TODO{3/3}\\
    \TODO{class 2} & \TODO{artifact} & \TODO{stage} & \TODO{3/3}\\
    \TODO{class 3} & \TODO{artifact} & \TODO{stage} & \TODO{3/3}\\
    \TODO{class 4} & \TODO{artifact} & \TODO{stage} & \TODO{3/3}\\
    \TODO{class 5} & \TODO{artifact} & \TODO{stage} & \TODO{3/3}\\
    \TODO{class 6} & \TODO{artifact} & \TODO{stage} & \TODO{3/3}\\
    \TODO{class 7} & \TODO{artifact} & \TODO{stage} & \TODO{3/3}\\
    \TODO{class 8} & \TODO{artifact} & \TODO{stage} & \TODO{3/3}\\
    \midrule
    total & & & 24/24\\
    \bottomrule
  \end{tabular}
  \caption{\TODO{PLACEHOLDER --- to be verified against the committed
  fault-injection results before submission.} The eight injected fault
  classes, the single artifact each mutates, the predeclared localization
  stage, and the masked evaluator's localization count across the three
  relation types.}
  \label{tab:faults}
\end{table}
